\Opensolutionfile{ans}[ans/ansTN-PhuBai111]

\noindent \textbf{PHẦN I. Câu trắc nghiệm một lựa chọn (3,0 đ)}
\vspace{0.2cm}

\begin{ex}	\macau{ID: VL11-PB111-C01}Phương trình dao động điều hòa của một chất điểm có dạng $x=A \cos (\omega t+\varphi)$. Độ dài quỹ đạo của dao động là
	\choice
	{$A$}
	{\True $2A$}
	{$4A$}
	{$A/2$}
	\loigiai{
		Quỹ đạo chuyển động trong không gian thực tế của một chất điểm dao động điều hòa là một đoạn thẳng giới hạn kéo dài từ biên âm ($-A$) đến biên dương ($+A$). Do đó, độ dài quỹ đạo chuyển động là $L = 2A$}
\end{ex}

\begin{ex}	\macau{ID: VL11-PB111-C02}Vật dao động điều hòa theo trục Ox. Phát biểu nào sau đây đúng?
	\choice
	{Lực kéo về tác dụng vào vật không đổi}
	{Quỹ đạo chuyển động của vật là một đường hình cos}
	{\True Quỹ đạo chuyển động của vật là một đoạn thẳng}
	{Li độ của vật tỉ lệ với thời gian dao động}
	\loigiai{
		\begin{itemize}
			\item Quỹ đạo của vật dao động điều hòa trong không gian là một đoạn thẳng có chiều dài $2A$. Đồ thị li độ dóng theo biến trục thời gian mới mang hình sin hoặc cosin. Do đó phương án C đúng.
			\item Lực kéo về biến thiên điều hòa theo thời gian ($F = -kx$), li độ biến thiên điều hòa lượng giác chứ không tỉ lệ tuyến tính đơn giản với thời gian.
		\end{itemize}
	}
\end{ex}

\begin{ex}	\macau{ID: VL11-PB111-C03}Dao động tắt dần
	\choice
	{\True có biên độ giảm dần theo thời gian}
	{luôn có lợi}
	{luôn có hại}
	{có biên độ không đổi theo thời gian}
	\loigiai{
		Theo định nghĩa chuẩn sách giáo khoa, dao động tắt dần là dao động cơ học có biên độ (và cơ năng toàn phần) bị giảm dần theo thời gian do tác dụng của lực cản và ma sát môi trường}
\end{ex}

\begin{ex}	\macau{ID: VL11-PB111-C04}Phát biểu nào dưới đây về dao động cưỡng bức là \textbf{sai}?
	\choice
	{Sau một thời gian dao động còn lại chỉ là dao động của ngoại lực tuần hoàn}
	{Tần số của dao động cưỡng bức bằng tần số của ngoại lực tuần hoàn}
	{Nếu ngoại lực cưỡng bức là tuần hoàn thì trong thời kì đầu dao động của con lắc là tổng hợp dao động riêng của nó với dao động của ngoại lực tuần hoàn}
	{\True Để trở thành dao động cưỡng bức, ta cần tác dụng lên con lắc dao động một ngoại lực không đổi}
	\loigiai{
		Để một hệ thực hiện dao động cưỡng bức ổn định, ngoại lực tác dụng bắt buộc phải là một ngoại lực biến thiên tuần hoàn theo thời gian (thường có dạng hàm sin hoặc cosin), không phải một lực không đổi. Nhận định D sai}
\end{ex}

\begin{ex}	\macau{ID: VL11-PB111-C05}Biên độ của hệ dao động điều hòa phụ thuộc vào yếu tố nào sau đây?
	\choice
	{Cách chọn gốc thời gian}
	{Cấu tạo của hệ}
	{\True Cách kích thích cho vật dao động}
	{Cách chọn trục tọa độ}
	\loigiai{
		Biên độ dao động điều hòa phụ thuộc hoàn toàn vào cách thức kích thích ban đầu cung cấp năng lượng cho hệ (ví dụ: kéo ra bao nhiêu rồi thả, hoặc truyền vận tốc bao nhiêu). Các đại lượng chu kì, tần số mới phụ thuộc cấu tạo của hệ}
\end{ex}

\begin{ex}	\macau{ID: VL11-PB111-C06}Đối với một vật dao động cưỡng bức giai đoạn ổn định:
	\choice
	{\True Chu kì dao động chỉ phụ thuộc vào ngoại lực}
	{Biên độ dao động chỉ phụ thuộc vào ngoại lực}
	{Chu kì dao động chỉ phụ thuộc vào vật và ngoại lực}
	{Biên độ dao động không phụ thuộc vào ngoại lực}
	\loigiai{
		Khi dao động cưỡng bức đã đi vào trạng thái ổn định lâu dài, chu kì và tần số dao động của hệ vật luôn bằng khít với chu kì và tần số của ngoại lực tuần hoàn kích thích, không phụ thuộc vào các thông số riêng của hệ vật}
\end{ex}

\begin{ex}	\macau{ID: VL11-PB111-C07}Dao động điều hòa có thể được coi như hình chiếu của một chuyển động tròn đều xuống một
	\choice
	{đường thẳng bất kì}
	{đường thẳng vuông góc với mặt phẳng quỹ đạo}
	{đường thẳng xiên góc với mặt phẳng quỹ đạo}
	{\True đường thẳng nằm trong mặt phẳng quỹ đạo}
	\loigiai{
		Hình chiếu vuông góc của một chất điểm chuyển động tròn đều xuống một đường thẳng (một trục tọa độ) nằm ngay bên trong mặt phẳng quỹ đạo tròn đó là một dao động điều hòa}
\end{ex}

\begin{ex}	\macau{ID: VL11-PB111-C08}Đối với một chất điểm dao động cơ điều hòa với chu kì T thì:
	\choice
	{Động năng và thế năng đều biến thiên tuần hoàn theo thời gian nhưng không điều hòa}
	{\True Động năng và thế năng đều biến thiên tuần hoàn điều hòa theo thời gian với chu kì bằng $T/2$}
	{Động năng và thế năng đều biến thiên tuần hoàn theo thời gian với chu kì $T$}
	{Động năng và thế năng đều biến thiên tuần hoàn theo thời gian với chu kì $2T$}
	\loigiai{
		Năng lượng động năng và thế năng của chất điểm dao động điều hòa tuần hoàn biến thiên điều hòa theo thời gian với tần số góc bằng $2\omega$, tần số bằng $2f$ và chu kì rút ngắn đi một nửa: $T' = \frac{T}{2}$}
\end{ex}

\begin{ex}	\macau{ID: VL11-PB111-C09}Trong dao động điều hòa có dạng phương trình $x=2A \cos (\omega t+\varphi)$, giá trị cực đại của gia tốc chất điểm là
	\choice
	{$a_{\max }=\omega^{2}A$}
	{\True $a_{\max }=2\omega^{2}A$}
	{$a_{\max }=2\omega^{2}A^{2}$}
	{$a_{\max }=-\omega^{2}A$}
	\loigiai{
		Giá trị cực đại của gia tốc luôn bằng tích số giữa bình phương tần số góc và biên độ dao động của vật. Vì biên độ đứng trước hàm của phương trình này có độ lớn bằng $2A$ nên gia tốc cực đại là: $a_{\text{max}} = 2\omega^2 A$}
\end{ex}

\begin{ex}	\macau{ID: VL11-PB111-C10}Gia tốc đại số của một vật dao động điều hòa sẽ triệt tiêu bằng không khi vật đi qua vị trí nào?
	\choice
	{\True Vật ở vị trí có li độ bằng không (VTCB)}
	{Vật ở vị trí có pha dao động cực đại}
	{Vật ở vị trí có li độ cực đại}
	{Vận tốc của vật đạt giá trị cực tiểu}
	\loigiai{
		Biểu thức định luật liên hệ gia tốc đại số: $a = -\omega^2 x$. Do đó, gia tốc bằng không ($a=0$) khi và chỉ khi li độ của vật triệt tiêu bằng không ($x=0$), tức là lúc vật đi qua vị trí cân bằng}
\end{ex}

\begin{ex}	\macau{ID: VL11-PB111-C11}In lòng dđđh của con lắc đơn, cơ năng toàn phần của hệ được xác định theo biên độ góc $\alpha_{0}$ (đơn vị rad), khối lượng $m$ của vật nặng và chiều dài $\ell$ của sợi dây theo công thức nào?
	\choice
	{$W=mg\ell\alpha_{0}^{2}$}
	{\True $W=\frac{1}{2}mg\ell\alpha_{0}^{2}$}
	{$W=\frac{1}{2}mg\alpha_{0}^{2}$}
	{$W=\frac{mg}{2\ell}\alpha_{0}^{2}$}
	\loigiai{
		Năng lượng cơ năng toàn phần của con lắc đơn thực hiện dao động điều hòa góc nhỏ ($\alpha_0 \le 10^{\circ}$) tuân theo hệ thức chuẩn: $W = \frac{1}{2}mg\ell\alpha_0^2$}
\end{ex}

\begin{ex}	\macau{ID: VL11-PB111-C12}Con lắc lò xo có độ cứng $k=20\text{ N/m}$ dao động điều hoà với biên độ $4\text{ cm}$. Động năng của con lắc khi đi qua vị trí có li độ $x=3\text{ cm}$ có giá trị bằng
	\choice
	{$0,1\text{ J}$}
	{$0,014\text{ J}$}
	{\True $0,007\text{ J}$}
	{$0,07\text{ J}$}
	\loigiai{
		\begin{itemize}
			\item Đổi đơn vị thông số kích thước hình học sang mét chuẩn: $A = 0,04\text{ m}$ và $x = 0,03\text{ m}$.
			\item Áp dụng định luật bảo toàn năng lượng, động năng bằng hiệu giữa cơ năng và thế năng:
			$$W_{\text{đ}} = W - W_{\text{t}} = \frac{1}{2}k(A^2 - x^2) = \frac{1}{2} \cdot 20 \cdot (0,04^2 - 0,03^2) = 10 \cdot 0,0007 = 0,007\text{ J} \text{}$$
		\end{itemize}}
\end{ex}

\Closesolutionfile{ans}

\noindent \textbf{PHẦN II. Câu trắc nghiệm đúng sai (2,0đ)}

\noindent \textit{Trong mỗi ý a), b), c), d) ở mỗi câu, thí sinh chọn đúng hoặc sai}
\Opensolutionfile{ans}[ans/ansTF-PhuBai111]

\begin{ex}	\macau{ID: VL11-PB111-C13}Một vật dao động điều hòa trên trục Ox quanh vị trí cân bằng O. Hình bên là đồ thị biểu diễn sự phụ thuộc của li độ x vào thời gian t.
% TODO: \usepackage{graphicx} required
\begin{center}
		\begin{tikzpicture}[>=stealth, line join=round, line cap=round, font=\footnotesize, xscale=1.2, yscale=0.5]
		% Định nghĩa các điểm tọa độ
		\path (0,0) coordinate (O);
		\path (0,3) coordinate (P0);
		\path (2,-6) coordinate (P1);
		\path (5,6) coordinate (P2);
		\path (8,-6) coordinate (P3);
		\path (2,0) coordinate (T1);
		\path (5,0) coordinate (T2);
		% Vẽ hệ trục tọa độ
		\draw[->] (-0.8, 0) -- (10, 0) node[below] {$t\text{ (s)}$};
		\draw[->] (0, -7.5) -- (0, 7.5) node[left] {$x\text{ (cm)}$};
		% Vẽ các đường gióng (dashed)
		\draw[dashed] (0, 6) -- (P2) -- (T2);
		\draw[dashed] (T1) -- (P1);
		\draw[dashed] (0, -6) -- (P3);
		% Vẽ đồ thị dao động
		\draw[thick] plot[domain=0:9.5, samples=150] (\x, {6*cos(deg(pi/3*\x + pi/3))});
		% Nhãn trên các trục
		\node[below left] at (O) {$0$};
		\node[left] at (0, 6) {$6$};
		\node[left] at (0, -6) {$-6$};
		\node[above] at (T1) {$\frac{1}{15}$};
		\node[below] at (T2) {$\frac{1}{6}$};
		% Đánh dấu các điểm đặc biệt
		\foreach \p in {P0, P1, P2, P3} {
			\fill[black] (\p) circle (1.5pt);
		}
	\end{tikzpicture}
\end{center}
	\choiceTF[t]
	{\True Biên độ dao động cực đại của vật có trị số bằng $6\text{ cm}$}
	{\True Tần số góc dao động riêng của hệ vật này có giá trị bằng $10\pi\text{ rad/s}$}
	{Tại thời điểm sau $t = \frac{1}{6}\text{ s}$ kể từ lúc bắt đầu chuyển động, vận tốc của vật đạt giá trị cực đại}
	{Vật đi qua vị trí có li độ $x = 1\text{ cm}$ tổng cộng 3 lần trong khoảng thời gian bằng $1,5T$ kể từ thời điểm ban đầu}
	\loigiai{
		\begin{itemchoice}
			\itemch Biên đỉnh cao nhất của đồ thị trục đứng cho biết biên độ dao động bằng $A = 6\text{ cm}$.
			\itemch Chu kì hình sóng trọn vẹn bám trục hoành chỉ ra $T = 0,2\text{ s} \Rightarrow \omega = \frac{2\pi}{T} = \frac{2\pi}{0,2} = 10\pi\text{ rad/s}$.
			\itemch Tại $t=0$, vật ở vị trí cân bằng và đi xuống theo chiều âm $\Rightarrow \varphi = \frac{\pi}{2}\text{ rad} \Rightarrow x = 6\cos(10\pi t + \frac{\pi}{2})$. Phương trình vận tốc: $v = -60\pi\sin(10\pi t + \frac{\pi}{2})$. Tại $t = \frac{1}{6}\text{ s} \Rightarrow v = -60\pi\sin(\frac{5\pi}{3} + \frac{\pi}{2}) = -60\pi\sin(\frac{13\pi}{6}) = -30\pi\text{ cm/s}$ (không đạt giá trị cực đại).
			\itemch Trong 1 chu kì $T$, vật qua vị trí $x=1\text{ cm}$ 2 lần. Trong nửa chu kì $0,5T$ tiếp theo, ta cần quét vòng tròn pha xem trạng thái quét cung có vấp phải điểm $x=1\text{ cm}$ lần thứ 4 hay không để đếm chính xác số lần.
		\end{itemchoice}
	}
\end{ex}

\begin{ex}	\macau{ID: VL11-PB111-C14}Hai vật dao động điều hòa cùng phương cùng tần số. Đồ thị biểu diễn động năng (đối với vật 1) hoặc thế năng (đối với vật 2) biến thiên phụ thuộc theo biến li độ x được mô tả như hình vẽ bên.
	% TODO: \usepackage{graphicx} required
% TODO: \usepackage{graphicx} required
\begin{center}
		\begin{tikzpicture}[>=stealth, line join=round, line cap=round, font=\footnotesize, scale=0.8]
		% Định nghĩa điểm
		\path (0,0) coordinate (O);
		% Vẽ lưới tọa độ
		\draw[step=1, gray!40, very thin] (-6, 0) grid (6, 4);
		% Vẽ trục tọa độ
		\draw[->, thick] (-7, 0) -- (7, 0) node[below] {$\lambda$};
		\draw[->, thick] (0, 0) -- (0, 5) node[right] {$W$};
		% Vẽ đồ thị (1) ĐN: y = 4 - 0.25*x^2
		\draw[thick, domain=-4:4, samples=100] plot (\x, {4 - 0.25*\x*\x});
		% Vẽ đồ thị (2) TN: y = 1/9*x^2
		\draw[thick, domain=-6:6, samples=100] plot (\x, {1/9*\x*\x});
		% Nhãn đồ thị và đường chỉ dẫn
		\draw[<-] (1.5, 3.4375) -- (2.5, 4.2) node[right] {(1) ĐN};
		\draw[<-] (-5, 2.7778) -- (-4, 4.2) node[above] {(2) TN};
		% Nhãn điểm
		\foreach \p/\g in {O/-135}
		\fill[black] (\p) circle (1.5pt) ++(\g:0.25) node{$\p$};
	\end{tikzpicture}
\end{center}
	
	\choiceTF[t]
	{Đường parabol có đỉnh hướng lên trên của vật (1) mô tả sự biến thiên của động năng phụ thuộc theo li độ}
	{\True Cơ năng toàn phần bảo toàn của hai hệ vật dao động này mang giá trị hoàn toàn bằng nhau}
	{Tỉ số biên độ dao động cực đại giữa vật (2) so với vật (1) mang giá trị bằng $2$}
	{\True Nếu tần số góc bằng nhau, tỉ số khối lượng của vật (1) gấp 2,25 lần khối lượng của vật (2)}
	\loigiai{
		\begin{itemchoice}
			\itemch Biểu thức động năng $W_{\text{đ}} = W - \frac{1}{2}kx^2$ là hàm bậc hai đối với $x$ có hệ số âm đứng trước biến, đồ thị là parabol có bề lõm quay xuống dưới, đỉnh hướng lên trên, nên mô tả đúng động năng [thành phần cắt trục tung tại giá trị lớn nhất vĩ mô].
			\itemch Điểm đỉnh cao nhất chạm thềm trục dọc của cả hai đồ thị đều gặp nhau chung tại một vạch mức trên trục năng lượng, chỉ ra cơ năng chúng bằng nhau: $W_1 = W_2 = W$.			\itemch Biên giới hạn hành trình biên trên trục hoành của đồ thị chỉ ra vật (1) có biên độ $A_1 = 4$ vạch ô lưới, vật (2) có biên độ $A_2 = 6$ vạch ô lưới. Tỉ số biên độ vật 2 so với vật 1 là: $\frac{A_2}{A_1} = \frac{6}{4} = 1,5$.
			\itemch Do cơ năng bằng nhau và tần số góc bằng nhau: $\frac{1}{2}m_1\omega^2 A_1^2 = \frac{1}{2}m_2\omega^2 A_2^2 \Rightarrow m_1 A_1^2 = m_2 A_2^2 \Rightarrow \frac{m_1}{m_2} = \left(\frac{A_2}{A_1}\right)^2 = (1,5)^2 = 2,25$.
		\end{itemchoice}
	}
\end{ex}

\Closesolutionfile{ans}

%%%%%%%%%%%%%%%%%%%%%%%%%%%%%%%%%%%%%%%%%%%%%%%%%%%%%%%%%%%%%%%%%%%%%%%%%
% PHẦN III: CÂU HỎI TRẮC NGHIỆM TRẢ LỜI NGẮN
%%%%%%%%%%%%%%%%%%%%%%%%%%%%%%%%%%%%%%%%%%%%%%%%%%%%%%%%%%%%%%%%%%%%%%%%%
\noindent \textbf{PHẦN III. CÂU HỎI NGẮN (2,0 điểm)}
\Opensolutionfile{ans}[ans/ansSA-PhuBai111]

\begin{ex}	\macau{ID: VL11-PB111-C15}Một vật dao động điều hoà có phương trình dao động dạng lượng giác: $x = 4 \cos\left(3t + \frac{\pi}{6}\right)\text{ (cm)}$. Tốc độ cực đại của vật khi đi qua vị trí cân bằng bằng bao nhiêu $\text{cm/s}$?

	\par\shortans[oly]{12}
	\loigiai{
		\begin{itemize}
			\item Trích xuất thông số biên độ $A = 4\text{ cm}$ và tần số góc $\omega = 3\text{ rad/s}$.
			\item Tốc độ cực đại của vật đạt được khi đi qua vị trí cân bằng bằng:
			$$v_{\text{max}} = \omega \cdot A = 3 \cdot 4 = 12\text{ cm/s} \text{}$$
		\end{itemize}
	}
\end{ex}

\begin{ex}	\macau{ID: VL11-PB111-C16}Một con lắc lò xo có khối lượng quả nặng $100\text{ g}$ thực hiện dao động cưỡng bức ổn định dưới tác dụng của ngoại lực biến thiên tuần hoàn theo thời gian với tần số $f$. Đồ thị biểu diễn sự phụ thuộc của biên độ dao động cưỡng bức theo tần số $f$ của ngoại lực được biểu diễn như hình vẽ bên. Chiếu đỉnh đồ thị cộng hưởng xuống trục hoành thu được vạch tần số bằng $2,5\text{ Hz}$. Hãy xác định độ cứng $k$ của lò xo theo đơn vị $\text{N/m}$ (lấy hằng số $\pi^2 = 10$).
	% TODO: \usepackage{graphicx} required
% TODO: \usepackage{graphicx} required
\begin{center}
		\begin{tikzpicture}[>=stealth, line join=round, line cap=round, font=\footnotesize, scale=1.2]
		% Trục tọa độ
		\draw[->, thick] (0, 0) -- (7.0, 0) node[below] {$\omega\text{ (rad/s)}$};
		\draw[->, thick] (0, 0) -- (0, 5.2) node[left] {$A\text{ (cm)}$};
		% Đường gióng phụ
		\draw[dashed] (0, 1.5) -- (6.0, 1.5);
		\draw[dashed] (1.5, 0) -- (1.5, 1.5);
		\draw[dashed] (6.0, 0) -- (6.0, 1.5);
		\draw[dashed] (0, 4.5) -- (3.75, 4.5) -- (3.75, 0);
		% Đồ thị dao động cưỡng bức (đường Gauss chuẩn hóa)
		\draw[thick, samples=100] plot[domain=0.8:6.7] (\x, {4.5*exp(-0.217*(\x-3.75)^2)});
		% Tọa độ các điểm nhãn
		\path (0,0) coordinate (O);
		\path (3.75,4.5) coordinate (F);
		% Nhãn trên các trục tọa độ
		\node[left] at (0, 1.5) {$4$};
		\node[left] at (0, 4.5) {$12$};
		\node[below] at (1.5, 0) {$2\pi$};
		\node[below] at (3.75, 0) {$5\pi$};
		\node[below] at (6.0, 0) {$8\pi$};
		% Nhãn điểm bằng vòng lặp ở cuối cùng
		\foreach \p/\g/\l in {O/-135/O, F/90/f_0}
		\fill[black] (\p) circle (1.5pt) ++(\g:0.25) node {$\l$};
	\end{tikzpicture}
\end{center}

	\par\shortans[oly]{25}
	\loigiai{
		\begin{itemize}
			\item Đỉnh của đồ thị cộng hưởng ứng với giá trị tần số xảy ra hiện tượng cộng hưởng cơ, tại đó tần số ngoại lực bằng tần số dao động riêng của hệ: $f_0 = 2,5\text{ Hz}$.
			\item Hệ thức tính tần số riêng của con lắc lò xo với $m = 100\text{ g} = 0,1\text{ kg}$:
			$$f_0 = \frac{1}{2\pi}\sqrt{\frac{k}{m}} \Rightarrow f_0^2 = \frac{k}{4\pi^2 m} \Rightarrow (2,5)^2 = \frac{k}{4 \cdot 10 \cdot 0,1}$$
			$$6,25 = \frac{k}{4} \Rightarrow k = 6,25 \cdot 4 = 25\text{ N/m} \text{}$$
		\end{itemize}
	}
\end{ex}

\begin{ex}	\macau{ID: VL11-PB111-C17}Một vật có khối lượng $200\text{ g}$, dao động điều hòa quanh vị trí cân bằng. Đồ thị hình bên mô tả động năng của vật ($W_{\text{đ}}$) thay đổi phụ thuộc vào thời gian $t$. Tại thời điểm ban đầu $t=0$, vật đang có li độ âm. Lấy hằng số $\pi^{2}=10$. Trích xuất dữ liệu chu kì biến thiên năng lượng từ hệ ô lưới đồ thị cho biết $T_{\text{nl}} = 0,2\text{ s}$ và đỉnh động năng cực đại đạt $W = 50\text{ mJ}$. Biên độ dao động dài của vật bằng bao nhiêu cm?
	% TODO: \usepackage{graphicx} required
% TODO: \usepackage{graphicx} required
\begin{center}
		\begin{tikzpicture}[line join=round, line cap=round, >=stealth, font=\footnotesize]
		% Grid
		\draw[dashed, gray!40, thin] (0, 0) grid[step=1] (6.5, 4.5);
		% Axes
		\draw[->, thick] (0, 0) -- (6.8, 0) node[below] {$t\text{ (s)}$};
		\draw[->, thick] (0, 0) -- (0, 4.8) node[left] {$W_{\text{đ}}\text{ (mJ)}$};
		% Axis labels
		\node[below left] at (0, 0) {$O$};
		\node[left] at (0, 2) {$20$};
		\node[left] at (0, 4) {$40$};
		\node[below] at (5, 0) {$0, 25$};
		% Curve
		\draw[thick, blue, smooth, samples=100, domain=0:6.3] plot (\x, {2 + 2*sin(\x*90)});
	\end{tikzpicture}
\end{center}

	\par\shortans[oly]{5}
	\loigiai{
		\begin{itemize}
			\item Đổi đơn vị khối lượng $m = 0,2\text{ kg}$ và động năng cực đại bằng cơ năng: $W = 50\text{ mJ} = 0,05\text{ J}$.
			\item Chu kì dao động của li độ gấp đôi chu kì biến thiên năng lượng: $T = 2 \cdot T_{\text{nl}} = 2 \cdot 0,2 = 0,4\text{ s}$.
			\item Tần số góc dao động riêng của hệ: $\omega = \frac{2\pi}{T} = \frac{2\pi}{0,4} = 5\pi\text{ rad/s}$.
			\item Sử dụng hệ thức liên hệ cơ năng toàn phần để tìm biên độ $A$:
			$$W = \frac{1}{2}m\omega^2 A^2 \Rightarrow 0,05 = \frac{1}{2} \cdot 0,2 \cdot (5\pi)^2 \cdot A^2 \Rightarrow 0,05 = 0,1 \cdot 25\pi^2 \cdot A^2$$
			Thế $\pi^2 = 10 \Rightarrow 0,05 = 25 \cdot A^2 \Rightarrow A^2 = 0,002 \Rightarrow A \approx 0,05\text{ m} = 5\text{ cm} \text{}$
		\end{itemize}
	}
\end{ex}

\begin{ex}	\macau{ID: VL11-PB111-C18}Một vật dao động điều hòa theo phương trình lượng giác: $x=4 \cos \left(\frac{2\pi}{3}t\right)$ ($x$ tính bằng cm, $t$ tính bằng s). Kể từ thời điểm ban đầu $t=0$, khoảng thời gian để vật đi qua vị trí có li độ $x=-2\text{ cm}$ lần thứ 2011 bằng bao nhiêu giây? Kết quả làm tròn lấy đến hàng số nguyên.

	\par\shortans[oly]{3016}
	\loigiai{
		\begin{itemize}
			\item Tần số góc $\omega = \frac{2\pi}{3}\text{ rad/s} \Rightarrow T = \frac{2\pi}{\omega} = 3\text{ s}$.
			\item Tại thời điểm $t=0$, vật ở biên dương ($x = +A = 4\text{ cm}$). Vị trí cần đi qua là $x = -2\text{ cm} = -\frac{A}{2}$.
			\item Trong một chu kì trọn vẹn $T$, chất điểm đi qua vị trí $-\frac{A}{2}$ đúng 2 lần.
			\item Tách số lần: $2011 = 2010 + 1 = 1005 \cdot 2 + 1$. 
			\item Khoảng thời gian hoàn thành 2010 lần đầu tương ứng mất hằng số hằng số thời gian bằng $1005T$. Sau đó vật quay trở về đúng biên dương.
			\item Để gặp vị trí $x = -\frac{A}{2}$ lần thứ 2011 (lần cuối), chất điểm đi tiếp từ biên dương đến vị trí $-\frac{A}{2}$ theo chiều âm. Góc quét tương ứng trên đường tròn lượng giác là $\alpha = \frac{\pi}{2} + \frac{\pi}{6} = \frac{2\pi}{3}\text{ rad}$. Khoảng thời gian đi đoạn này là $\Delta t = \frac{\alpha}{\omega} = \frac{\frac{2\pi}{3}}{\frac{2\pi}{3}} = 1\text{ s}$.
			\item Tổng thời gian toàn phần đạt:
			$$t = 1005 \cdot T + \Delta t = 1005 \cdot 3 + 1 = 3015 + 1 = 3016\text{ s} \text{}$$
		\end{itemize}
	}
\end{ex}

\Closesolutionfile{ans}

\noindent \textbf{PHẦN IV. TỰ LUẬN (3 điểm)}
\vspace{0.2cm}

\begin{ex}	\macau{ID: VL11-PB111-C19}Một vật dao động điều hòa có phương trình li độ dạng: $x=2 \cos \left(4 \pi t+\frac{\pi}{3}\right)(\mathrm{cm})$. Hãy xác định:
	\begin{enumerate}
		\item[a/] Chu kì riêng của dao động.
		\item[b/] Li độ đại số của vật tại thời điểm $t=2\text{ s}$.
	\end{enumerate}
	\loigiai{
		\begin{enumerate}
			\item[a/] Tần số góc từ phương trình là $\omega = 4\pi\text{ rad/s}$. Chu kì riêng của dao động:
			$$T = \frac{2\pi}{\omega} = \frac{2\pi}{4\pi} = 0,5\text{ s} \text{}$$
			\item[b/] Thay thời điểm giá trị $t=2\text{ s}$ trực tiếp vào phương trình li độ của vật:
			$$x(2) = 2\cos\left(4\pi \cdot 2 + \frac{\pi}{3}\right) = 2\cos\left(8\pi + \frac{\pi}{3}\right) = 2\cos\left(\frac{\pi}{3}\right) = 2 \cdot 0,5 = 1\text{ cm} \text{}$$
		\end{enumerate}
	}
\end{ex}

\begin{ex}	\macau{ID: VL11-PB111-C20}Một chất điểm có khối lượng $200\text{ g}$ dao động điều hòa với tần số riêng $f=1\text{ Hz}$. Tại thời điểm ban đầu mốc chọn $t=0$, vật đi qua vị trí có li độ $x=5\text{ cm}$ với tốc độ hiển thị $v=100\pi\sqrt{3}\text{ cm/s}$ (theo hướng dẫn giải chi tiết tương ứng $10\pi\sqrt{3}\text{ cm/s}$) đang chuyển động theo chiều dương.
% TODO: \usepackage{graphicx} required
\begin{center}
		\begin{tikzpicture}[>=stealth, line join=round, line cap=round, font=\footnotesize, scale=0.8]
		% Lưới tọa độ
		\draw[cyan!40!blue!30, thin] (-4, 0) grid (4, 8);
		% Trục tọa độ
		\draw[->, thick] (-4.8, 0) -- (5.2, 0) node[below] {$x\text{ (cm)}$};
		\draw[->, thick] (0, -0.5) -- (0, 9) node[right] {$W_{\text{đ}}\text{ (mJ)}$};
		% Đường biểu diễn (Parabol)
		\draw[very thick, blue!80!black, domain=-4:4, samples=100] plot (\x, {-0.5*\x*\x + 8});
		% Vạch chia và nhãn trên trục hoành
		\foreach \x in {-4, -2, 2, 4} {
			\draw[thick] (\x, 0.1) -- (\x, -0.1) node[below=2pt] {$\x$};
		}
		\node[below left] at (0, 0) {$0$};
		% Vạch chia và nhãn trên trục tung
		\foreach \y/\label in {2/20, 4/40, 6/60} {
			\draw[thick] (0.1, \y) -- (-0.1, \y) node[left=2pt] {\label};
		}
		\draw[thick] (0.1, 8) -- (-0.1, 8);
		\node[above right] at (0, 8) {$80$};
	\end{tikzpicture}
\end{center}
	\begin{enumerate}
		\item[a.] Tính giá trị biên độ dao động của vật?
		\item[b.] Viết phương trình vận tốc theo biến thời gian của dao động trên.
	\end{enumerate}
	\loigiai{
		\begin{enumerate}
			\item[a.] Tần số góc dao động riêng của chất điểm: $\omega = 2\pi f = 2\pi \cdot 1 = 2\pi\text{ rad/s}$.
			Sử dụng biểu thức công thức độc lập vuông pha để xác định trị số biên độ $A$:
			$$A = \sqrt{x^2 + \left(\frac{v}{\omega}\right)^2} = \sqrt{5^2 + \left(\frac{10\pi\sqrt{3}}{2\pi}\right)^2} = \sqrt{25 + (5\sqrt{3})^2} = \sqrt{25 + 75} = 10\text{ cm} \text{}$$
			\item[b.] Biên độ của hàm vận tốc đạt độ lớn cực đại: $v_{\text{max}} = \omega A = 2\pi \cdot 10 = 20\pi\text{ cm/s}$.
			Tại $t=0 \Rightarrow x_0 = A\cos\varphi \Rightarrow 5 = 10\cos\varphi \Rightarrow \cos\varphi = 0,5 \Rightarrow \varphi = \pm \frac{\pi}{3}\text{ rad}$. Vì vật đi theo chiều dương thuận nên chọn pha ban đầu li độ dấu âm: $\varphi = -\frac{\pi}{3}\text{ rad}$.
			Pha ban đầu của vận tốc sớm pha hơn li độ một góc vuông: $\varphi_v = \varphi + \frac{\pi}{2} = -\frac{\pi}{3} + \frac{\pi}{2} = \frac{\pi}{6}\text{ rad}$.
			Phương trình vận tốc theo thời gian của dao động là:
			$$v = 20\pi\cos\left(2\pi t + \frac{\pi}{6}\right)\text{ cm/s} \text{}$$
		\end{enumerate}
	}
\end{ex}

\begin{ex}	\macau{ID: VL11-PB111-C21}Đồ thị hình dưới đây mô tả sự thay đổi của hàm động năng $W_{\text{đ}}$ theo li độ $x$ của quả cầu khối lượng $m$ bên trong một con lắc lò xo dao động điều hòa theo phương thẳng đứng. Đỉnh cao nhất đồ thị dóng sang trục tung cho biết vạch cơ năng đạt $0,08\text{ J}$ và biên hoành độ đạt biên độ $A = 4\text{ cm}$. Tính:
	\begin{enumerate}
		\item[a.] Độ cứng $k$ của lò xo.
		\item[b.] Giá trị động năng của con lắc lò xo khi quả cầu đi qua vị trí có li độ $x = 2\text{ cm}$.
	\end{enumerate}
	\loigiai{
		\begin{enumerate}
			\item[a.] Đổi đơn vị kích thước biên độ sang hệ mét chuẩn SI: $A = 4\text{ cm} = 0,04\text{ m}$. Đỉnh cao nhất của động năng bằng giá trị hằng số cơ năng toàn phần của hệ: $W = 0,08\text{ J}$.
			Áp dụng công thức cơ năng tính độ cứng $k$ của lò xo:
			$$W = \frac{1}{2}kA^2 \Rightarrow 0,08 = \frac{1}{2} \cdot k \cdot (0,04)^2 \Rightarrow 0,08 = 0,0008k \Rightarrow k = 100\text{ N/m} \text{}$$
			\item[b.] Khi quả cầu đi qua vị trí có li độ $x = 2\text{ cm} = 0,02\text{ m}$, thế năng đàn hồi của hệ bằng:
			$$W_{\text{t}} = \frac{1}{2}kx^2 = \frac{1}{2} \cdot 100 \cdot (0,02)^2 = 50 \cdot 0,0004 = 0,02\text{ J} \text{}$$
			Theo định luật bảo toàn, giá trị động năng tương ứng của hệ bằng:
			$$W_{\text{đ}} = W - W_{\text{t}} = 0,08 - 0,02 = 0,06\text{ J} \text{}$$
		\end{enumerate}
	}
\end{ex}

\Closesolutionfile{ans}

\newpage
%%%%%%%%%%%%%%%%%%%%%%%%%%%%%%%%%%%%%%%%%%%%%%%%%%%%%%%%%%%%%%%%%%%%%%%%%
% KHUNG ĐÁP ÁN BẢNG TỰ ĐỘNG IN KẾT XUẤT CỦA HỆ THỐNG
%%%%%%%%%%%%%%%%%%%%%%%%%%%%%%%%%%%%%%%%%%%%%%%%%%%%%%%%%%%%%%%%%%%%%%%%%
\begin{center} \textbf{BẢNG ĐÁP ÁN THAM KHẢO LỚP 11 - TRƯỜNG PHÚ BÀI} \end{center}

\noindent \textbf{Phần I (Câu hỏi trắc nghiệm một phương án lựa chọn):}